% Part: computability
% Chapter: computability-theory
% Section: coding-computations

\documentclass[../../../include/open-logic-section]{subfiles}

\begin{document}

\olfileid[hi]{cmp}{thy}{cod}
\olsection{संगणनाओं का कूटांकन}

संगणना के प्रत्येक मॉडल में निम्नलिखित करना संभव है:
\begin{enumerate}
\item संगणनीय फलनों की \emph{परिभाषाओं} का व्यवस्थित ढंग से वर्णन करना।
  उदाहरणतः, ट्यूरिंग मशीन के विनिर्देशों, पुनरावर्ती परिभाषाओं अथवा किसी
  प्रोग्रामिंग भाषा के प्रोग्रामों को ऐसी परिभाषाएँ देने वाला माना जा सकता है।
\item किसी दी हुई परिभाषा वाले फलन की किसी दी हुई आगत पर संगणना का पूरा
  अभिलेख वर्णित करना। उदाहरणतः, किसी ट्यूरिंग मशीन की संगणना को संगणना के
  प्रत्येक चरण के विन्यासों (मशीन की अवस्था और टेप की सामग्री) के अनुक्रम
  द्वारा वर्णित किया जा सकता है।
\item यह जाँचना कि संगणना का कोई प्रस्तावित अभिलेख वास्तव में इस बात का
  अभिलेख है कि दी हुई परिभाषा वाले संगणनीय फलन को किसी दी हुई आगत पर कैसे
  संगणित किया जाएगा (जिस पर फलन परिभाषित है, अर्थात् संगणना रुकती है)।
\item संगणना के पूर्ण अभिलेख के ऐसे वर्णन से किसी दी हुई आगत के लिए फलन का
  मान निकालना। उदाहरणतः, रुकने वाली ट्यूरिंग मशीन की संगणना के अंतिम चरण में
  टेप की सामग्री ही वह मान होती है।
\end{enumerate}

कूटांकन के द्वारा किसी संगणनीय फलन के प्रत्येक वर्णन को एक संख्यात्मक
\emph{सूचकांक} इस प्रकार दिया जा सकता है कि सूचकांक से निर्देश संगणनीय ढंग
से पुनः प्राप्त किए जा सकें। इसी तरह, किसी संगणना के पूरे अभिलेख को भी एक
ही संख्या से कूटित किया जा सकता है। तब इससे प्राप्त अंकगणितीय संबंध
``$s$, सूचकांक~$e$ वाले फलन की आगत~$x$ पर संगणना का अभिलेख कूटित करता
है'' और फलन ``कूट~$s$ वाले संगणना-अनुक्रम का निर्गत'' संगणनीय होते हैं;
वास्तव में वे आदिम पुनरावर्ती होते हैं।

यह मूलभूत तथ्य अत्यंत शक्तिशाली है और चुने गए संगणना-मॉडल से स्वतंत्र
रहकर संगणनीयता के विषय में अनेक उल्लेखनीय और महत्त्वपूर्ण परिणाम सिद्ध
करने देता है।

\end{document}
